\documentclass[11pt]{article}
\usepackage[margin=0.9in]{geometry}
\usepackage[T1]{fontenc}
\usepackage[utf8]{inputenc}
\usepackage[hidelinks]{hyperref}
\usepackage{xurl}
\usepackage{graphicx}
\usepackage{enumitem}
\usepackage{amsmath}
\usepackage{tikz}
\usetikzlibrary{arrows.meta, positioning, shapes.geometric}
\setlist[itemize]{topsep=0.25em, itemsep=0.25em}
\setlist[enumerate]{topsep=0.25em, itemsep=0.25em}

\newcommand{\projectname}{Quantyze}
\newcommand{\projecttitle}{Trees, Trades, and Books}
\newcommand{\projectsubtitle}{Simulating a Limit Order Book with Tree-Based Data Structures and an ML Agent}
\newcommand{\projectdate}{March 31, 2026}

\begin{document}
\begin{center}
{\Huge \projectname: \projecttitle\par}
\vspace{0.8em}
\noindent\rule{0.82\textwidth}{0.6pt}\par
\vspace{0.9em}
{\large\itshape\parbox{0.9\textwidth}{\centering \projectsubtitle}\par}
\vspace{1.4em}
\begin{minipage}[t]{0.31\textwidth}
\centering
\large Cade McNelly
\end{minipage}
\hfill
\begin{minipage}[t]{0.31\textwidth}
\centering
\large Nicolas Miranda Cantanhede
\end{minipage}
\hfill
\begin{minipage}[t]{0.31\textwidth}
\centering
\large Sahand Samadirand
\end{minipage}

\vspace{1.2em}
{\large \projectdate\par}
\end{center}
\vspace{1.2em}

\section*{1. Introduction}
As a group of CS students interested in finance, we have spent a lot of time
looking at stock charts. It does not take long to notice a major disconnect
between the price shown on the screen and the mechanics operating underneath
it. A market interface may display a stock at a certain level, but that does
not mean an order can be filled there instantly or in unlimited size. In
practice, fills depend on liquidity, queue position, and the slippage created
when an order moves through the book.

That gap between a quoted price and an executable price is one of the main
ideas that motivated this project. It is also a very natural place to apply
CSC111 ideas such as tree-based data structures, simulation, and algorithmic
state updates. Quantyze studies that gap by simulating a limit order book in
pure Python. We represent bids and asks as separate binary search trees of
price levels, maintain FIFO priority within each level, replay event streams
through a matching engine, and compute summary statistics about fills,
cancellations, spread, midpoint, and slippage. On top of the simulator, we
added a small machine learning component that observes book-state
features and classifies market actions in the short term.

\textbf{Our project goal is to build a tree-based market simulator that can
replay realistic order flow, expose the mechanics of price-time-priority
matching, and train a small neural network to classify short-horizon market
actions as buy, sell, or hold from order-book state.}

The project kept the same core direction as our proposal, but the final version
organizes the submission around a reliable browser-first workflow launched from
\texttt{main.py}, with a terminal fallback for the same simulator and training
logic. That let us present the project in a PyCharm-friendly way without
changing the main computational challenge: representing and updating a live
order book efficiently with trees while connecting the simulator to a
meaningful data and model-training pipeline. In the language of our original
proposal, we still see the project as an attempt to close the gap between what
a trader sees on a chart and what actually has to happen for an order to be
filled.

\section*{2. Datasets}
One thing that became much clearer while building the final project is that not
all data in Quantyze plays the same role. Quantyze uses three distinct
\textbf{data-source roles}, and separating them clearly makes the rest of the
system much easier to follow:
\begin{itemize}
    \item \textbf{synthetic scenarios} for simulator-first replay demos
    \item \textbf{packaged internal CSVs} for fast and reproducible classifier
    retraining
    \item \textbf{paired raw LOBSTER files} for the real-data preprocessing and
    feature-building path
\end{itemize}

The submission also includes a packaged dataset zip, so the TA does not need to
download anything externally before running the project. That was important to
us because we wanted the grading path to be self-contained and easy to repeat.

\subsection*{2.1 Synthetic Event Data}
We generated our own synthetic event streams for testing and demonstration.
These are created inside \texttt{data\_loader.py} and include three scenarios:
\begin{itemize}
    \item \texttt{balanced}
    \item \texttt{low\_liquidity}
    \item \texttt{high\_volatility}
\end{itemize}

These synthetic datasets use the same canonical event format as the rest of the
project: \texttt{timestamp}, \texttt{order\_id}, \texttt{side},
\texttt{order\_type}, \texttt{price}, and \texttt{quantity}. Their role is
\textbf{simulation}, not packaged retraining. They are the fastest way for the
TA to see fills, cancellations, spread changes, and book updates without
waiting for a large dataset to load.

\subsection*{2.2 LOBSTER Sample Data}
For the real-data path of the machine learning portion of the project, we used
LOBSTER sample order-book data from the official LOBSTER sample page.\footnote{\url{https://data.lobsterdata.com/info/DataSamples.php}}
Each LOBSTER sample comes as a \textbf{pair of files}:
\begin{itemize}
    \item a message file containing timestamped market events
    \item an order-book file containing depth snapshots
\end{itemize}

Our code supports raw LOBSTER message and order-book CSVs. The training flow is
pointed at the \textbf{message file}, and the loader infers the paired
order-book file automatically from the filename. From the message file, we use
normalized information corresponding to timestamp, event type, order id, size,
price, and direction. The loader also recognizes message types that should be
preserved as annotations rather than replayed directly.

From the paired order-book file, our current feature pipeline uses only the top
two bid levels and top two ask levels, even when a file contains deeper book
data. In practice, the model ultimately trains on level-1 and level-2
price/size information rather than the full book.

During development, we used public LOBSTER sample files such as AAPL, AMZN, and
GOOG to validate the real-data pipeline. In the final submission, we ship one
official Apple 5-level sample pair inside the dataset zip using the clearer
extracted filenames:
\begin{itemize}
    \item \texttt{aapl\_lobster\_2012-06-21\_message\_5level\_sample.csv}
    \item \texttt{aapl\_lobster\_2012-06-21\_orderbook\_5level\_sample.csv}
\end{itemize}

This packaged LOBSTER option is \textbf{optional}. It demonstrates the
real-data pipeline, but it is much larger and slower to load than the short
internal packaged dataset.

\subsection*{2.3 Submission Dataset Package}
For the final submission, the dataset zip is \texttt{quantyze\_datasets.zip}. It
contains:
\begin{itemize}
    \item \texttt{sample\_internal.csv}
    \item \texttt{huge\_internal.csv}
    \item \texttt{aapl\_lobster\_2012-06-21\_message\_5level\_sample.csv}
    \item \texttt{aapl\_lobster\_2012-06-21\_orderbook\_5level\_sample.csv}
    \item \texttt{instructions.txt}
    \item \texttt{model.pt}
    \item \texttt{training\_metrics.json}
\end{itemize}

The intended roles of these packaged files are:
\begin{itemize}
    \item \texttt{sample\_internal.csv}: the \textbf{quick} retraining check
    \item \texttt{huge\_internal.csv}: the \textbf{large internal} retraining
    demo, which takes longer to load and train
    \item \texttt{aapl\_lobster\_2012-06-21\_message\_5level\_sample.csv}: the
    \textbf{large raw LOBSTER message-file} demo
    \item \texttt{aapl\_lobster\_2012-06-21\_orderbook\_5level\_sample.csv}:
    the paired order-book file required by the raw LOBSTER path
\end{itemize}

The packaged saved-state workflow also includes a shipped baseline checkpoint
and metrics artifact, so the TA can use the model path without retraining
first.

Because this zip is under the course upload limit, we include
\texttt{quantyze\_datasets.zip} directly in the final MarkUs submission rather
than asking the TA to download it from an external service.

\section*{3. Computational Overview}
\subsection*{3.1 Core Data Representation}
The central data structure in Quantyze is the limit order book. We represent the
book with two binary search trees:
\begin{itemize}
    \item one tree for bids
    \item one tree for asks
\end{itemize}

This logic is implemented across \texttt{book\_tree.py},
\texttt{price\_level.py}, and \texttt{order\_book.py}. Each node in a tree is a
\texttt{PriceLevel} object. A price level stores:
\begin{itemize}
    \item a price
    \item aggregate resting volume at that price
    \item a FIFO queue of resting orders
\end{itemize}

The trees let us preserve sorted price structure while still supporting
efficient best-bid and best-ask access. The FIFO queue inside each price level
enforces time priority among orders resting at the same price. This combination
directly models price-time-priority matching, which was the main algorithmic
goal of our project.

\subsection*{3.2 Event Replay and Matching}
All incoming data is converted into a common \texttt{Event} representation in
\texttt{orders.py}. The \texttt{DataLoader} in \texttt{data\_loader.py} either:
\begin{itemize}
    \item loads a dataset from internal CSV or raw LOBSTER files, or
    \item generates a synthetic scenario
\end{itemize}

Those events are then replayed through \texttt{EventStream} in
\texttt{event\_stream.py}, which emits them one at a time into the
\texttt{MatchingEngine} in \texttt{matching\_engine.py}. The matching engine
handles limit orders, market orders, and cancellations. It updates the
\texttt{OrderBook}, performs partial and full fills, removes empty price levels
when necessary, and records execution information used for later metrics.

At the end of a normal run, \texttt{main.py} prints a terminal summary with:
\begin{itemize}
    \item total filled quantity
    \item fill count
    \item cancel count
    \item average slippage
    \item final spread
    \item final midpoint
    \item agent overlay mark-to-market P\&L if a trained checkpoint is active
\end{itemize}

The current default run path is command-line based rather than browser based.
The project therefore reports results through terminal output, the packaged
saved-state artifacts extracted from the dataset zip, and runtime logs such as:
\begin{itemize}
    \item \texttt{model.pt}
    \item \texttt{training\_metrics.json}
    \item \texttt{log.json}
\end{itemize}

\subsection*{3.3 Data Processing and Training Pipeline}
The machine learning pipeline is built mainly in \texttt{data\_loader.py},
\texttt{neural\_net.py}, and \texttt{main.py}.

The feature pipeline uses 12 features per example, extracted by
\texttt{DataLoader.\_feature\_vector\_from\_levels}:
\begin{enumerate}
    \item best bid price
    \item best bid size
    \item best ask price
    \item best ask size
    \item spread
    \item mid-price
    \item level-1 imbalance
    \item second bid price
    \item second bid size
    \item second ask price
    \item second ask size
    \item event-side feature
\end{enumerate}

The label policy is a \textbf{next-changed-mid rule}, implemented in
\texttt{DataLoader.\_labels\_from\_mid\_sequence}. For each example, we scan
forward through the mid-price sequence until the midpoint changes by a
non-zero amount:
\begin{itemize}
    \item if the next observed change is an increase, the label is
    \texttt{buy}
    \item if the next observed change is a decrease, the label is
    \texttt{sell}
    \item if no change is observed within the remaining sequence, the label
    is \texttt{hold}
\end{itemize}

This rule labels each snapshot with the direction of the very next midpoint
move, rather than comparing against a fixed future horizon. It avoids an
arbitrary look-ahead window and produces a label that is directly tied to
the next price-relevant event in the replay.

This pipeline is implemented in \texttt{DataLoader.build\_training\_dataset} and
related helper methods. The final training flow in \texttt{main.py} computes
train/validation splits, standardizes features with training-set mean and
standard deviation, applies class-weighted cross-entropy, trains a PyTorch
feed-forward network, and saves the latest training artifacts:
\begin{itemize}
    \item \texttt{latest\_model.pt}
    \item \texttt{latest\_training\_metrics.json}
    \item \texttt{latest\_training\_data.csv}
\end{itemize}

The shipped baseline saved-state artifacts remain \texttt{model.pt} and
\texttt{training\_metrics.json}, so the TA can always return to a known working
checkpoint.

\subsection*{3.4 User Interactivity}
The TA-facing interface is a nested terminal menu implemented in
\texttt{cli\_menu.py}. Running \texttt{main.py} with no arguments opens a menu
with separate sections for:
\begin{itemize}
    \item quick TA demo
    \item simulation
    \item training
    \item artifacts and metrics
    \item help / command reference
\end{itemize}

The quick TA demo is the recommended first path because it prints the current
overlay status, runs the default \texttt{balanced} scenario, and explains how
to interpret the resulting metrics. The rest of the menu lets the TA choose
scenarios, replay datasets, run short training, inspect baseline versus latest
metrics, inspect artifact status, inspect the current execution log, inspect
the packaged dataset zip, and switch the simulation overlay between the shipped
baseline checkpoint, the latest retrained checkpoint, and no model. The active
choice is stored in
\texttt{active\_model.json}, which makes the training-to-simulation workflow
persistent across runs. In the final submission, the browser UI and the
terminal fallback share the same runtime entry point in \texttt{main.py}.

\subsection*{3.5 Overall Computational Complexity}
Quantyze demonstrates computational complexity in all three major ways
highlighted by the project rubric.

\begin{itemize}
    \item \textbf{Non-trivial dataset processing and model training:} the system
    supports multiple input formats, infers paired LOBSTER order-book files,
    separates replay-compatible and non-replayable events, builds a shared
    training dataset, normalizes features, computes class weights, and saves
    checkpoints and metric artifacts.
    \item \textbf{Non-trivial algorithms implemented by us:} the BST-based order
    book, FIFO price-level queues, price-time-priority matching engine, and
    cancellation/indexing logic are all custom project code rather than thin
    wrappers around third-party libraries.
    \item \textbf{Non-trivial reporting and interactivity:} the browser UI,
    terminal fallback, saved-state workflow, metrics inspection, and
    execution-log reporting make the program substantially more interactive
    than a single fixed command.
\end{itemize}

\section*{4. Instructions for Obtaining Datasets and Running the Program}
The TA-facing run path is organized as a workflow rather than a list of
independent commands. The most important distinction is that
\textbf{simulation source} and \textbf{model source} are separate:
\begin{itemize}
    \item the \textbf{simulation source} is the event stream being replayed,
    such as \texttt{balanced}, \texttt{low\_liquidity},
    \texttt{high\_volatility}, or a custom CSV
    \item the \textbf{model source} is the classifier overlay currently
    selected for inference, such as the shipped \texttt{model.pt} baseline or
    a newly trained \texttt{latest\_model.pt}
\end{itemize}

Keeping them separate lets the TA test the same classifier overlay under
different market conditions and retrain a model without overwriting the shipped
baseline artifacts. We found this distinction important because the simulator
is the core system, while the model is an optional layer on top of it.

\subsection*{4.1 Installing Libraries}
Install the libraries listed in \texttt{requirements.txt} before running the
program.

\subsection*{4.2 Obtaining Datasets}
The TA downloads \texttt{quantyze\_datasets.zip} from our final MarkUs
submission and extracts it beside \texttt{main.py}. This provides:
\begin{itemize}
    \item a small internal CSV for a quick retraining check
    \item a larger internal CSV for a slower baseline-scale retraining demo
    \item a packaged Apple LOBSTER \textbf{message file} and its paired
    \textbf{order-book file}
    \item the baseline checkpoint and metrics artifacts
\end{itemize}

The shortest TA path should start with \texttt{sample\_internal.csv}. The
larger internal dataset and the packaged LOBSTER pair are both slower-loading
options for a more thorough optional check.

\subsection*{4.3 Workflow Overview}
The primary grading path is the browser UI launched by running
\texttt{main.py}. At each stage, it makes four things explicit:
\begin{itemize}
    \item What event source is being replayed?
    \item Which checkpoint, if any, is active?
    \item Where did that checkpoint come from?
    \item What is the next logical action: inspect, train, or simulate?
\end{itemize}

The intended workflow starts with the default browser run, then branches into
optional extra simulations, classifier retraining, or fallback terminal use
through \texttt{main.py --cli}.

\subsection*{4.4 Recommended TA Demo}
The clearest way to understand the project is the following walkthrough.

Open the project in PyCharm, run \texttt{main.py} with no additional
arguments, and open the printed local URL in a browser.

We recommend the following sequence:
\begin{enumerate}
    \item Start on the \texttt{Simulate} tab and run the default
    \texttt{balanced} scenario. This is the fastest simulator-first path and
    produces the same summary metrics described later in this section.
    \item On the same tab, optionally run \texttt{low\_liquidity} and
    \texttt{high\_volatility}. These scenarios show that the same overlay can
    be evaluated under different event streams without changing the core
    matching engine.
    \item Open \texttt{Train}. For a short demonstration, use
    \texttt{sample\_internal.csv}. For a baseline-scale demonstration, use
    \texttt{huge\_internal.csv}. \\
    For the real-data path, use
    \texttt{aapl\_lobster\_2012-06-21\_message\_5level\_sample.csv}. The
    latter two are both much larger and take longer to load. After training,
    the UI writes the latest output files and exposes the new checkpoint in the
    saved artifact status.
    \item Open \texttt{Artifacts}. Check the overlay status, saved files, and
    baseline versus latest metrics. The provenance should now point to the
    newly trained checkpoint rather than the shipped baseline if the latest
    model has been activated.
    \item Open \texttt{Charts} after a simulation. This shows order-book depth
    and execution-price history based on the replay that just ran.
    \item Re-run the \texttt{balanced} synthetic demo under a different
    overlay. This uses the same event
    stream as before, but under a different overlay. It helps isolate the
    difference between \textit{market replay} and \textit{agent overlay}
    behavior.
\end{enumerate}

\subsection*{4.5 How the TA Should Interpret the Output}
The browser UI and direct terminal output contain the same two classes of
information, and keeping them separate is the key to understanding Quantyze.

\subsubsection*{Simulation Run Configuration}
Before each run, the program prints:
\begin{itemize}
    \item the \textbf{event source}, such as \texttt{synthetic balanced}
    or a custom replay path
    \item the \textbf{simulation overlay mode}, such as \texttt{baseline},
    \texttt{latest}, or \texttt{none}
    \item the \textbf{simulation overlay path}
    \item the \textbf{overlay provenance}, such as the packaged baseline
    checkpoint or a newly trained dataset like \texttt{sample\_internal.csv}
\end{itemize}

This header tells the TA which market is being replayed and which classifier
overlay, if any, is observing it. They are related, but they are not the same
thing.

\subsubsection*{Matching-Engine Metrics}
The following summary values describe the replay itself:
\begin{itemize}
    \item \texttt{Total Filled}: total executed quantity produced by the
    matching engine
    \item \texttt{Fill Count}: number of executions
    \item \texttt{Cancel Count}: number of cancellations processed
    \item \texttt{Average Slippage}: average execution-price deviation under the
    current slippage calculation
    \item \texttt{Spread}: final bid-ask spread
    \item \texttt{Mid Price}: final midpoint
\end{itemize}

These values are properties of the event stream and the matching engine. If the
TA replays the same scenario with \texttt{baseline}, \texttt{latest}, and
\texttt{none}, the market metrics should stay the same because the underlying
event stream is unchanged.

\subsubsection*{Agent Overlay Metric}
If an active checkpoint is loaded, the program also prints
\texttt{Agent Overlay Mark-to-Market P\&L}. This value needs to be interpreted
carefully.

It is \textbf{not} a claim that the model is a profitable trading strategy. In
the current system, the neural network is trained to classify the direction of the next
non-zero midpoint change: buy, sell, or hold. During
simulation, that prediction is converted into a simple paper-trading overlay
that increments or decrements inventory by one unit whenever the simulator
observes fills. As a result:
\begin{itemize}
    \item a strong classification accuracy does not automatically imply a strong
    trading P\&L
    \item a negative mark-to-market value can reflect inventory accumulation in
    a thin market rather than a broken simulator
    \item the ML component should be interpreted as a proof-of-concept
    prediction overlay, not as a finished trading strategy
\end{itemize}

For example, \texttt{low\_liquidity} often ends with a wider spread and fewer
fills than \texttt{balanced}. A model can therefore finish with non-zero
inventory and a negative mark-to-market value even when the matching engine
itself is behaving correctly.

\subsubsection*{Training Metrics}
The training menu prints validation accuracy, majority-class baseline,
per-class recall, and the confusion matrix. These values measure how well the
classifier reproduces the label rule derived from future midpoint movement.
They do not directly measure whether the simulated agent will earn positive
P\&L on every synthetic scenario.

\subsection*{4.6 Optional Raw LOBSTER Check}
If the TA wants to test the raw LOBSTER path, they do \textbf{not} need to
download external files. They can use the packaged training-menu option for
\texttt{aapl\_lobster\_2012-06-21\_message\_5level\_sample.csv}. The paired
order-book path is inferred automatically from that message filename, which
maps to
\texttt{aapl\_lobster\_2012-06-21\_orderbook\_5level\_sample.csv}.

The same loader logic also works for a custom raw LOBSTER message CSV path: it
infers the paired order-book path by replacing \texttt{\_message\_} with
\texttt{\_orderbook\_}.

\section*{5. Changes from Proposal to Final Submission}
Several parts of the final project changed in useful ways from the proposal.

The simulator core became stronger and more concrete. The final system has a
clearer separation between loading data, replaying events, matching orders,
computing summaries, and handling both the browser UI and the terminal
fallback. The BST-based order
book and FIFO price-level logic were central from the beginning, but the final
project makes their interaction much clearer.

The machine learning plan also became more specific. In the proposal, we left
room for either a decision-tree or neural-network predictor. In the final
submission, we committed to a PyTorch neural network and built a full training
pipeline with explicit features, normalization, class weighting, saved
checkpoints, and validation metrics. In other words, the ML side became less
open-ended and more concrete.

Our display plan also became more realistic. The proposal emphasized a more
ambitious React/Plotly-style interface, but the final submission focuses on a
lighter browser UI launched directly from \texttt{main.py} and backed by the
same Python simulator and training code as the CLI fallback. We chose to
prioritize a robust tree-based simulator and a reproducible training workflow
over a more ambitious frontend stack. That was a practical choice, not a
change in what we thought the project was about.

The data plan also became more realistic. Instead of treating external data in
the abstract, the final project clearly separates synthetic scenarios for
replay, packaged internal CSVs for fast retraining, and a packaged raw LOBSTER
pair for the real-data preprocessing path. That keeps the TA workflow
self-contained while still showing that the project can process real market
data.

\section*{6. Discussion}
Our goal was to model the mechanics of a limit order book using trees and to
connect that simulator to a simple market-action classifier. The final project
does meet that goal, though the machine learning side is more modest than the
simulator itself.

On the simulation side, the output is meaningful and easy to interpret. In the
current default \texttt{balanced} run, the simulator reports 1748.0 total
filled volume, 374 fills, 128 cancellations, a final spread of approximately
0.6, and a final midpoint of 100.0. Those values show that the event sequence
is not just inserting passive orders forever; it is generating executions,
consuming liquidity, and rebuilding the book over time. That mattered to us
from the beginning, because one of the main motivations for the project was to
show the difference between a quoted market price and the mechanics required to
fill an order.

We also think the project succeeds as a data-structures project. The binary
search trees are not decorative. They are the central representation of the
market. Maintaining sorted price levels, extracting the best bid and best ask,
removing empty levels, and preserving FIFO priority within a level all fit the
tree-based design naturally. That gives the project real computational
substance beyond simple data loading or basic statistics.

The machine learning component works, but it needs to be interpreted carefully.
The packaged baseline checkpoint and its accompanying metrics artifact show a
validation accuracy of approximately 0.99992 compared to a majority-class
baseline of approximately 0.33753. Those numbers are useful for showing that
the training pipeline, checkpoint loader, and saved-state workflow function
correctly, but they should not be read as a realistic claim about
real-market prediction.

The run-time P\&L shown in the simulator should also be read as an agent
overlay metric rather than as proof of profitability. The neural network is
trained on future midpoint labels, not on a direct trading reward, and the
simulation overlay uses those labels to build a simple paper position over
time. Because of that, the classifier can behave as designed while still
producing a negative mark-to-market value in a thin or volatile synthetic
scenario.

When we tested on public LOBSTER sample data during development, the task was
much harder and performance was only modestly above a naive baseline. For that
reason, we do not present the current model as a strong predictive system for
real markets. In this project, its role is better understood as a small 
proof of concept that connects book-state features to model output inside a
complete simulator.

That still answers the project goal in a meaningful way. The goal was never
just to maximize predictive accuracy. It was to build a full computational
pipeline from market events to order-book state to derived features to model
output. In that sense, the ML side is successful: it uses both synthetic and
real-data paths, produces saved checkpoints, supports inference inside the
simulator, and can be retrained from the browser UI, the command line, or the
terminal fallback menu.

We also ran into clear limitations. Synthetic data is easy to test and reason
about, but it is not enough to justify a strong predictive claim. Real LOBSTER
data is much richer, but it brings preprocessing issues, class-balance issues,
and a harder learning problem. The model itself is also simple: we used a small
feed-forward network over handcrafted features rather than a stronger temporal
model. That was a reasonable CSC111 choice, but it limits what the model can
learn about market dynamics. Finally, our presentation scope narrowed over the
course of the project. The proposal imagined a more polished UI layer, but most
of our time went into the simulation engine, dataset handling, and ML pipeline.

If we were to continue the project, the next steps would be to improve the
visualization layer, add a stronger sequence-aware model, and evaluate the
trained agent with more realistic trading metrics across multiple market days.
Even with those limitations, we think the current submission succeeds as a
CSC111 project because it combines tree-based design, non-trivial event
processing, model training, and interpretable output in one coherent system.
More importantly, it still reflects the original motivation behind the project:
to make the hidden mechanics of a market feel concrete rather than abstract.

\section*{7. References}
\begin{enumerate}
    \item LOBSTER. (n.d.). \textit{Sample data}. Retrieved March 31, 2026, from
    \url{https://data.lobsterdata.com/info/DataSamples.php}
    \item PyTorch. (n.d.). \textit{PyTorch}. Retrieved March 31, 2026, from
    \url{https://pytorch.org/}
    \item Investopedia. (n.d.). \textit{Limit order book}. Retrieved March 31,
    2026, from
    \url{https://www.investopedia.com/terms/l/limitorderbook.asp}
    \item Investopedia. (n.d.). \textit{Depth of market (DOM)}. Retrieved March
    31, 2026, from
    \url{https://www.investopedia.com/terms/d/depth-of-market.asp}
    \item QuestDB. (n.d.). \textit{Limit order book}. Retrieved March 31, 2026,
    from \url{https://questdb.com/glossary/limit-order-book/}
    \item University of Oxford. (2019). \textit{Introduction to limit order book
    markets}. Retrieved March 31, 2026, from
    \url{https://www.maths.ox.ac.uk/system/files/attachments/NUS-LOB19.pdf}
\end{enumerate}

\end{document}
